\begin{trapbox}
	\begin{description}[style=nextline, leftmargin=2.8em, labelsep=0.5em, itemsep=0.35em]
		\item[\textbf{\textcolor{CTrap}{易错点一（符号顺序错）：}}] 将公式错记为 $n(A)+n(B)+n(C) + n(A\cap B) + \cdots$。
		\item[\textbf{\textcolor{CTrap}{正确理解（加减交替）：}}] 正确顺序是先加单集，再减去两两交集，最后加回三重交集。
		\item[\textbf{\textcolor{CTrap}{错因分析（记忆混淆）：}}] 与二集合公式或不等式变形混淆，未理解“修正重复计数”的逻辑。
	\end{description}

	\medskip
	\begin{description}[style=nextline, leftmargin=2.8em, labelsep=0.5em, itemsep=0.35em]
		\item[\textbf{\textcolor{CTrap}{易错点二（遗漏三交）：}}] 仅减去两两交集而未加回三重交集，导致结果偏小。
		\item[\textbf{\textcolor{CTrap}{正确理解（必须加回）：}}] 三重交集在减去两两交时被多次扣除，必须加回一次才能准确。
		\item[\textbf{\textcolor{CTrap}{错因分析（算理不清）：}}] 机械记忆公式而忽略每一步的计数意义，以为减去所有两两交已足够。
	\end{description}

	\medskip
	\begin{description}[style=nextline, leftmargin=2.8em, labelsep=0.5em, itemsep=0.35em]
		\item[\textbf{\textcolor{CTrap}{易错点三（忽略有限前提）：}}] 对无限集或元素个数不明的集合强行使用该公式。
		\item[\textbf{\textcolor{CTrap}{正确理解（仅限有限集）：}}] $n(\cdot)$ 表示有限集的元素个数，无限集无此概念，公式无意义。
		\item[\textbf{\textcolor{CTrap}{错因分析（忽视条件）：}}] 看到集合符号就直接套用，未确认是否“有限”以及数据是否完备。
	\end{description}
\end{trapbox}
